\documentclass[11pt]{article}

\usepackage[margin=1in]{geometry} 
\usepackage{amsmath,amsthm,amssymb,enumitem}
\usepackage[sectionbib]{natbib}

\newcommand{\R}{\mathbb{R}}
\newcommand{\Z}{\mathbb{Z}}
\newcommand{\N}{\mathbb{N}}
\newcommand{\Q}{\mathbb{Q}}
\newcommand{\C}{\mathbb{C}}
\newcommand{\F}{{\mathbb F}}
\newcommand{\Pb}{{\mathbb P}}
\newcommand{\E}{{\mathbb E}}
\newcommand{\Frac}[2]{\left(\frac{#1}{#2}\right)}
\newcommand{\Eval}[2]{\big|_{#1}^{#2}}
\newcommand{\intinf}{\int_{-\infty}^{\infty}}

\newenvironment{theorem}[2][Theorem]{\begin{trivlist}
\item[\hskip \labelsep {\bfseries #1}\hskip \labelsep {\bfseries #2.}]}{\end{trivlist}}
\newenvironment{lemma}[2][Lemma]{\begin{trivlist}
\item[\hskip \labelsep {\bfseries #1}\hskip \labelsep {\bfseries #2.}]}{\end{trivlist}}
\newenvironment{exercise}[2][Exercise]{\begin{trivlist}
\item[\hskip \labelsep {\bfseries #1}\hskip \labelsep {\bfseries #2.}]}{\end{trivlist}}
\newenvironment{problem}[2][Problem]{\begin{trivlist}
\item[\hskip \labelsep {\bfseries #1}\hskip \labelsep {\bfseries #2.}]}{\end{trivlist}}
\newenvironment{question}[2][Question]{\begin{trivlist}
\item[\hskip \labelsep {\bfseries #1}\hskip \labelsep {\bfseries #2.}]}{\end{trivlist}}
\newenvironment{corollary}[2][Corollary]{\begin{trivlist}
\item[\hskip \labelsep {\bfseries #1}\hskip \labelsep {\bfseries #2.}]}{\end{trivlist}}

\newenvironment{solution}{\begin{proof}[Solution]}{\end{proof}}

\begin{document}

% ------------------------------------------ %
%                 START HERE                 %
% ------------------------------------------ %

\title{Assignment 1} % Replace X with the appropriate number
\author{ \iffalse STUDENT NAME \fi \\ECEV 33400} % Replace  "\iffalse X \fi" with your name.
\date{Due Date: 10/5/2017} % Don't change this! :( 

\maketitle

% -----------------------------------------------------
% The following two environments (theorem, proof) are
% where you will enter the statement and proof of your
% first problem for this assignment.
%
% In the theorem environment, you can replace the word
% "theorem" in the \begin and \end commands with
% "exercise", "problem", "lemma", etc., depending on
% what you are submitting. Replace the "x.yz" with the
% appropriate number for your problem.
%
% If your problem does not involve a formal proof, you
% can change the word "proof" in the \begin and \end
% commands with "solution".
% -----------------------------------------------------

\subsection*{Deterministic {\em vs} Stochastic Dynamics:\\ A theory of island biogeography in the limit of $S \rightarrow \infty$}
As you know, in 1967, MacArthur \& Wilson published a book ``The Theory of Island Biogeography'' (TIB) to explain the factors controlling species diversity on {\em islands} at certain distance from the continent, the {\em mainland}. The stochastic model underlying this theory will be reviewed around the 7th week. A simpler version of it makes the following assumptions:
\begin{enumerate}
  \item {\bf $S \rightarrow \infty$:} The number of species, $S$, in the mainland is very large. Technically, $S$ has tended to $\infty$.
  \item {\bf Species equivalence:} Species arrive in an island at the same rate, $\lambda$, the colonization rate, and, they get extinct at the same rate, $\mu$, as well.
  \item {\bf Species independence:} Species arrivals and extinctions are random independent events. 
\end{enumerate}

Given the previous assumptions, address the following questions:
\begin{enumerate}
\item Suppose you have an island with a certain number of species $n$. What is the deterministic equation governing the temporal evolution of this number? Use this equation to calculate the number of species, $n^{\star}$, at the colonization-extinction equilibrium? 

\item The probability a new species arrives in time $dt$ is given by $\Pb[n \rightarrow n+1] = \lambda\ dt$. Since all species become extinct at the same rate, $\mu$, independently, the probability of loosing any of them to extinction is given by $\Pb[n \rightarrow n -1] = \mu\ n\ dt$. Use this probabilities to build the set of differential equations for the temporal evolution of the probability distribution describing $n$. This equation is the stochastic counterpart of the previous deterministic equation. 

\item Calculate the stationary distribution for the number of species at equilibrium, $\Pb^{\star}(n)$. What is the expected value for this number?

\item Consider a system of $K$ islands undergoing the colonization-extinction process described by the three assumptions above.  Assume there is no inter-island colonization so that all islands receive species directly from the mainland at rate $\lambda$. Assume also the system has reached equilibrium. If species on an island are detected with probability $a$, what is the probability of observing $n_i$ species on the $i$-th island? 

\item Assume you visited the archipelago of $K$ islands for a survey and observed $n_1, \dots, n_K$ species on the 1st to the $K$th islands, respectively. Assume again that there is no inter-island colonization so that all islands receive species directly from the mainland at rate $\lambda$. The sample of species observed with detectability $a$ across islands can be represented by the vector $\vec{n} = (n_1, \dots, n_K)$. Assume the process underlying this sample is described by the three assumptions above and that the whole archipelago is at colonization-extinction equilibrium. Calculate the probability of obtaining the sample $\vec{n}$, $\Pb^{\star}(\vec{n}) | \lambda, \mu)$. 

\item The probability $\Pb^{\star}(\vec{n}) | \lambda, \mu)$ is sometimes called a sampling probability distribution. It can be regarded also as the likelihood of obtaining the sample $\vec{n}$ given that model parameters are $\lambda$, $\mu$, and $a$. If we define 
\begin{equation}
  F(a, \lambda, \mu \,|\, \vec{n}) \equiv \Pb^{\star}(\vec{n} \,|\, a, \lambda, \mu) 
\end{equation}
and take logarithms, we obtain: 
\begin{equation}
  \log \left( F(a, \lambda, \mu \,|\, \vec{n} ) \right) = \log \left( \Pb^{\star}(\vec{n} \,|\, a, \lambda, \mu) \right) 
% * <dalonso@ceab.csic.es> 2017-10-17T15:31:45.487Z:
%
% ^.
\end{equation}
which is called the loglikelihood function. By defining the auxiliary variable $x \equiv a \lambda/\mu$, calculate the parameter $x^{\star}$ that maximizes the probability of obtaining the sample vector $\vec{n}$. 
\end{enumerate}

% -----------------------------------------------------
% First problem
% -----------------------------------------------------

% \begin{problem}{1}
% Show that the probability-generating function for the binomial distribution $p_j = \binom{n}{j} p^j q^{n-j}$, where $ q = 1 - p$, is $P(x) = (q + px)^n$. Hence prove that the mean is $m = np$, and the variance is $\sigma^2 = npq$.
% \end{problem}

% \begin{solution}
% \end{solution}

% -----------------------------------------------------
% Second problem
% -----------------------------------------------------

% \vspace{0.25in} % This just adds some space between problems 1 and 2.

% \begin{problem}{2}
% What is the probability-generating function for the Poisson distribution $$p_j = \frac{e^{-\lambda} \lambda^j}{j!}, j = 0, 1, 2, \cdots ?$$ Hence show that the mean and variance of the distribution are both equal to $\lambda$. 
% \end{problem}

% \begin{solution}
% \end{solution}

% -----------------------------------------------------
% Third problem
% -----------------------------------------------------


% -----------------------------------------------------
% Fourth problem
% -----------------------------------------------------


% -----------------------------------------------
% Ignore everything that appears below this.
% -----------------------------------------------

\end{document}

\documentclass[11pt]{article}
\usepackage{amssymb,amsmath,amsthm,mathrsfs,enumitem}

%%%%%%%%%%%%%%%%
%Change these values to the correct ones for your class
\newcommand{\CourseName}{STAT 25100}
\newcommand{\HWNum}{3}
\newcommand{\StudentName}{Ruby An}
\newcommand{\InstructorName}{Professor Ding}
\newcommand{\DueDate}{\today}
%%%%%%%%%%$%%%%%

%You can define commands for the things you use frequently.
\newcommand{\Z}{{\mathbb Z}}
\newcommand{\R}{{\mathbb R}}
\newcommand{\N}{{\mathbb N}}
\newcommand{\F}{{\mathbb F}}
\newcommand{\Pb}{{\mathbb P}}
\newcommand{\E}{{\mathbb E}}
\newcommand{\Frac}[2]{\left(\frac{#1}{#2}\right)}
\newcommand{\Eval}[2]{\big|_{#1}^{#2}}
\newcommand{\intinf}{\int_{-\infty}^{\infty}}


%We define environments for questions and solutions (so in the body of the document you can type \begin{question} ... \end{question} and \begin{solution} ... \end{solut
\newenvironment{question}[1]{\noindent{\large{\bf{Problem #1.}}} \begin{em}}{\end{em}\bigskip}
\newenvironment{solution}{\noindent\ignorespaces\par\noindent \large{\bf{Solution}} }{\bigskip}

%Set up the margins
\usepackage[margin=1in]{geometry}

%Create a header at the top of every page
\usepackage{fancyhdr}
\pagestyle{fancy}
\fancyhead[L]{\CourseName \\ \InstructorName} %Put the name on the Left 
\fancyhead[C]{Assignment \HWNum}  %Put the Course name on the Center 
\fancyhead[R]{\StudentName \\ \DueDate}    %Put the Assignment Number on the Right 

%Start the document
\begin{document}

\begin{question}{1}
\end{question}

\begin{solution}\\
\end{solution}


\begin{question}{2}

\begin{enumerate}[label=\bf{(\alph*)}]
\item
\end{enumerate}
\end{question}

\begin{solution}\\

\begin{enumerate}[label=\bf{(\alph*)}]
\item 
\end{enumerate}
\end{solution}



\end{document}

%Question Template
\begin{question}

\end{question}
\begin{solution}

\end{solution}

\begin{enumerate}[label=\bf{(\alph*)}]
\item
\end{enumerate}



